\begin{problem}[第32届全国中学生物理竞赛决赛]{53}
太阳系内行星的公转方向是基本相同的。地球上的观测者所看到的行星位置实际上是该行星在遥远的背景星空中的投影。由于各行星的公转速度以及它们在轨道上位置的不同，地球上的观测者看到的行星在背景星空中移动的方向与地球公转方向并非总是相同。人们把看到的行星移动方向与地球公转方向相同时的行星视运动叫顺行，方向相反的叫逆行。当行星位于由顺行转成逆行或由逆行转成顺行的转折点时，行星看来好像停留在星空不动，该位置叫留。天文学把地外行星运行到太阳和地球所在的直线上、且太阳和地外行星位于地球的两侧的情形叫行星冲日。火星冲日是常见的天文现象。设地球和火星都在同一平面内绕太阳做圆周运动，且只考虑太阳对行星的引力。已知火星轨道半径 $R_{\mathrm{m}}$ 为地球轨道半径 $R_{0}$ 的1.524倍，不考虑地球自转，试计算火星在经历相邻两次冲日的时间间隔内其视运动为逆行和顺行的时间间隔。
\end{problem}
